% ============================================================
%  EXTENDED ENGLISH SUMMARY
%  Hostifer — A Native Kubernetes PaaS Platform for Algeria
%  Master's Thesis — Computer Science, USTHB
% ============================================================

\chapter*{Extended Summary in English}
\addcontentsline{toc}{chapter}{Extended Summary in English}

\begin{center}
    {\large\textbf{Hostifer: Design and Implementation of a Native\\
    Kubernetes Multi-Tenant PaaS Platform for the Algerian Market}}\\[12pt]
    {\normalsize Master's Thesis in Computer Science\\
    University of Science and Technology Houari Boumediene (USTHB)\\
    Faculty of Computer Science}
\end{center}

\vspace{16pt}

% ── 1. Introduction ──────────────────────────────────────────
\section*{1. Introduction and Motivation}

The rapid growth of the Algerian digital ecosystem — with over 36 million internet users as of early 2025, more than 7,800 registered startups, and a 60\% year-over-year increase in startup investment to 650 million USD in 2024 — has created a pressing demand for developer infrastructure that local actors can access, afford, and trust. However, the dominant Platform-as-a-Service (PaaS) offerings available on the global market present a set of structural barriers that make them effectively inaccessible or legally non-compliant for Algerian developers and organisations.

International platforms such as Vercel, Railway, Render, Heroku, Fly.io, and Koyeb require payment in USD via international credit cards, a barrier that excludes the majority of Algerian freelancers and early-stage startups whose banking access is limited to DZD-denominated cards. None of these platforms operate data centres on Algerian territory, meaning that any personal data processed through their infrastructure is subject to cross-border transfer — a practice that requires prior authorisation from the National Authority for the Protection of Personal Data (ANPDP) under Law 18-07. The enforcement of Law 25-11 in July 2025, which strengthens data residency obligations and introduces new accountability requirements including a mandatory Data Protection Officer and a five-day breach notification deadline, has made this compliance gap increasingly consequential.

Existing Algerian cloud initiatives fail to address these needs adequately. Hawiyat focuses primarily on VPS hosting and no-code automation tools rather than automated deployment pipelines. Deployi offers traditional hosting services without any containerisation abstraction. PivoCloud provides container deployment based on Docker Compose without the orchestration, isolation, or autoscaling capabilities required for production multi-tenant workloads. Sagittario is non-functional at the time of writing. Only OpenScaler presents a technically serious IaaS/CaaS offering, though it remains in public alpha and targets infrastructure provisioning rather than the developer-facing PaaS experience.

Hostifer is designed to fill this gap precisely. It is a cloud-native PaaS platform, built on Kubernetes, hosted on Algerian infrastructure, priced in Algerian Dinars, and architecturally compliant with Laws 18-07 and 25-11. It enables any developer to deploy a web application from a GitHub repository URL in under five minutes, without writing a Dockerfile, configuring a server, or managing infrastructure primitives. This thesis documents the design decisions, architecture, and implementation of the platform across its three main components: the NestJS API backend, the Go build worker, and the Kubernetes infrastructure layer.

% ── 2. State of the Art ──────────────────────────────────────
\section*{2. State of the Art}

\subsection*{2.1 Cloud Service Models}

The National Institute of Standards and Technology (NIST) defines three fundamental cloud service models: Infrastructure as a Service (IaaS), Platform as a Service (PaaS), and Software as a Service (SaaS), differentiated by the level of abstraction offered to the consumer and the distribution of operational responsibility between provider and user. IaaS provides raw compute, storage, and networking resources, leaving operating system and application management to the consumer. SaaS delivers fully managed applications over the internet with no user control over the underlying stack. PaaS occupies the middle tier: the provider manages all infrastructure, and the consumer retains control only over the application code, its configuration, and the deployment pipeline. Hostifer is unambiguously a PaaS platform — developers interact exclusively with their source code and the platform's web interface, with no exposure to the Kubernetes primitives operating beneath.

\subsection*{2.2 Container Orchestration with Docker and Kubernetes}

Docker, launched in 2013, popularised containerisation as a lightweight virtualisation model in which an application and its complete dependency stack are packaged into a portable, reproducible OCI-compliant image. The Open Container Initiative (OCI), established in 2015, standardised the image format and runtime specification, enabling interoperability across tools. Kubernetes, originally developed at Google and donated to the Cloud Native Computing Foundation (CNCF) in 2016, emerged as the de facto standard for orchestrating containerised workloads at scale. Its architecture separates a control plane (API Server, Scheduler, Controller Manager, etcd) from a data plane of worker nodes executing container workloads via a CRI-compliant runtime such as containerd. Seven Kubernetes resource types are central to the implementation of Hostifer: Pod (the atomic execution unit), Namespace (logical cluster partition), Deployment (replica lifecycle manager), Service (stable network endpoint), Ingress (external HTTP routing), ResourceQuota (namespace resource cap), and NetworkPolicy (network traffic firewall).

\subsection*{2.3 International PaaS Platforms}

A comparative analysis of six leading international PaaS platforms reveals consistent patterns of inaccessibility for the Algerian context. Vercel, optimised for frontend and serverless workloads built on Next.js, bills exclusively in USD and operates no infrastructure in Africa. Railway, the creator of the Railpack build tool used in Hostifer, offers usage-based pricing inaccessible without an international credit card. Render provides predictable monthly instance pricing but no African deployment region. Heroku, the pioneer of the PaaS paradigm, eliminated its free tier in 2022 and now occupies a premium price tier relative to its competitors. Fly.io deploys applications as Firecracker microVMs across 18+ global regions but has no African presence and requires CLI proficiency that represents a significant barrier for non-infrastructure engineers. Koyeb offers a serverless edge platform with global reach but plans starting at 79 USD/month, far beyond the budget of most Algerian startups. None of these platforms provide a mechanism for data residency compliance under Algerian law.

\subsection*{2.4 The Algerian Regulatory Framework}

Law 18-07 of 10 June 2018 on the protection of personal data established Algeria's primary data protection framework, introducing obligations of informed consent, security safeguards, breach notification, and — critically — a prohibition on cross-border data transfers without prior authorisation from the ANPDP. Law 25-11 of 24 July 2025 strengthens and extends this framework, introducing mandatory Data Protection Officer designation, an automated processing registry, a five-day breach notification deadline, and an explicit principle of data sovereignty requiring that personal data of Algerian residents be stored and processed on Algerian territory. Any Algerian organisation using a foreign cloud provider to host applications that collect personal data of Algerian users is therefore potentially non-compliant with both laws, exposing itself to administrative penalties under ANPDP enforcement. The growing legislative pressure toward data residency creates a structural market demand for a compliant local PaaS offering that existing infrastructure providers cannot satisfy.

% ── 3. System Design ─────────────────────────────────────────
\section*{3. System Design and Architecture}

\subsection*{3.1 Actors and Requirements}

The system was designed around five categories of actors: the \textbf{Developer} (the primary human actor, interacting via the Next.js frontend to deploy and manage applications), the \textbf{Platform Administrator} (a secondary human actor supervising infrastructure), \textbf{GitHub} (an external automated actor triggering CI/CD pipelines via webhook events), \textbf{Argo Workflows} (the internal orchestration engine executing deployment pipelines), and the \textbf{Kubernetes API} (the infrastructure layer processing all resource provisioning requests). Supporting automated actors include NextAuth.js, the Go Builder Job, Railpack CLI, Buildkitd, Helm, HashiCorp Vault with the External Secrets Operator, Traefik, cert-manager, and the Promtail-Loki observability chain.

Functional requirements were prioritised using the MoSCoW method across seven domains: authentication (GitHub OAuth and email/password with JWT rotation), project management (creation from GitHub URL, settings, deletion with full deprovisioning), deployment (automated build, provision, and DNS pipeline triggered manually or via webhook), real-time log streaming (Server-Sent Events), environment variable management (Vault-backed secret storage), CI/CD configuration (automatic deployment on push), and platform administration. Non-functional requirements set measurable targets: end-to-end deployment time under 5 minutes (BNF-01), SSE log latency below 2 seconds (BNF-02), 99.5\% availability SLA for Standard plans (BNF-03), complete inter-tenant network isolation validated by NetworkPolicy (BNF-05), TLS encryption on all user-facing communications (BNF-07), and multi-layer autoscaling via HPA, VPA, and Cluster Autoscaler (BNF-09 through BNF-11).

\subsection*{3.2 UML Modelling}

The system was modelled using four categories of UML diagrams. Use case diagrams document three primary user interactions: the multi-step Deploy Wizard (covering GitHub account connection, repository selection, deployment configuration, build log streaming, and live URL access), environment variable management (including Vault-backed CRUD operations and secret reveal with audit logging), and project deletion (covering parallel Helm uninstall, DNS record deletion, and OCI registry image removal). Sequence diagrams detail the authentication flow with OAuth delegation and refresh token rotation, the full deployment pipeline from frontend submission to LIVE state, the SSE log streaming chain from pod stdout through Promtail, Loki, and the NestJS SSE endpoint to the browser, and the GitHub webhook CI/CD trigger flow. The class diagram defines eleven entities in the PostgreSQL schema: User, Project, Deployment, DeploymentStep, EnvVar, Plan, GitHubInstallation, GitHubRepository, ActivityLog, Notification, and Domain. A component diagram maps the layered architecture across eight logical layers: Client, API, Data, Orchestration, Build, Infrastructure, Security, Networking, and Observability.

\subsection*{3.3 Technical Architecture}

The platform architecture separates concerns across five distinct repositories: \texttt{frontend} (Next.js TypeScript), \texttt{backend} (NestJS TypeScript), \texttt{helm} (main application Helm chart), \texttt{tenant} (per-tenant Kubernetes provisioning chart), and \texttt{build-job} (Go binary). Three Go microservice repositories handle DNS lifecycle: \texttt{create-dns-job}, \texttt{delete-dns-job}. All compiled artefacts are published to the GitHub Container Registry (GHCR) as OCI images under the \texttt{ghcr.io/hostifer/} namespace. Infrastructure state is managed via a \texttt{gitops} repository consumed by ArgoCD for GitOps-based continuous delivery.

The platform is hosted on Alibaba Cloud Container Service for Kubernetes (ACK Managed Pro) in Algeria. The cluster comprises two node pools: a control pool of three ECS instances (4 vCPU / 16 GB RAM) running platform system workloads, and a worker pool of three ECS instances (4 vCPU / 8 GB RAM) running tenant application workloads, horizontally extensible via the ACK Cluster Autoscaler. Network architecture employs two Server Load Balancers (ports 80/443 for application traffic, port 6443 for Kubernetes API access), three Elastic IP Addresses (two inbound, one outbound associated with an Internet NAT Gateway for SNAT), and a dedicated VPC with a tenant node vSwitch. Storage is tiered across NAS (NFS, ReadWriteMany) for the OCI registry and Loki, ESSD cloud disks for Buildkitd layer cache, and OSS object storage for cold log archival.

Technological choices were driven by three criteria: functional fit, maturity, and stack coherence. NestJS was selected over Express and Fastify for its modular architecture, native Passport integration, and SSE decorator support. Prisma was chosen over TypeORM and Drizzle for its declarative schema, type-safe generated client, and deterministic migrations. Argo Workflows was preferred over Tekton and Jenkins X for its Kubernetes-native WorkflowTemplate model, built-in UI, and fine-grained pod lifecycle policies. Railpack and Buildkit were selected jointly over Kaniko and Cloud Native Buildpacks for their combined ability to detect frameworks automatically without requiring a Dockerfile and to produce OCI images with parallel layer caching.

% ── 4. Implementation ─────────────────────────────────────────
\section*{4. Implementation}

\subsection*{4.1 NestJS Backend}

The NestJS backend is structured as seventeen independent modules, each encapsulating a distinct functional domain. The \texttt{auth/} module implements three authentication mechanisms: email/password authentication using bcrypt hashing and Passport JWT strategy, GitHub OAuth delegation via \texttt{passport-github2} with complete flow handled server-side to prevent token exposure in the browser, and two-factor authentication via TOTP. Refresh token rotation is implemented in \texttt{refreshTokens()}: each use of a refresh token invalidates it immediately and issues a new pair, with refresh tokens stored as bcrypt hashes in PostgreSQL to prevent exposure in case of database compromise. The NextAuth.js JWT callback on the frontend detects access token expiry and calls \texttt{POST /auth/refresh} automatically, making session renewal transparent to the user.

The \texttt{deployments/} module is responsible for the full deployment lifecycle. Upon receiving a \texttt{POST /deployments} request, it creates a \texttt{Deployment} record in status \texttt{QUEUED}, then submits the \texttt{hostifer-deploy} WorkflowTemplate to Argo Workflows via its REST API with five dynamic parameters: \texttt{repo\_url}, \texttt{subdomain}, \texttt{port}, \texttt{tenant\_id}, and \texttt{image\_name}. Status synchronisation is performed via periodic polling of the Argo REST endpoint \texttt{GET /api/v1/workflows/\{namespace\}/\{name\}}, mapping Argo phases (\texttt{Pending}, \texttt{Running}, \texttt{Succeeded}, \texttt{Failed}) to internal \texttt{DeploymentStatus} values (\texttt{QUEUED}, \texttt{BUILDING}/\texttt{PROVISIONING}/\texttt{CONFIGURING\_DNS}, \texttt{LIVE}, \texttt{FAILED}). Each status transition updates both the \texttt{Deployment} record and the corresponding \texttt{DeploymentStep} records with timing metadata.

Real-time build log streaming is implemented at \texttt{GET /deployments/:id/logs} using the NestJS \texttt{@Sse()} decorator, which returns an RxJS \texttt{Observable<MessageEvent>} consumed by the browser's native \texttt{EventSource} API. The underlying async generator polls Loki's \texttt{/loki/api/v1/query\_range} endpoint every two seconds, filtering by the LogQL stream selector \texttt{\{build\_id="<deploymentId>", container="main"\}}. The \texttt{build\_id} label is attached by Promtail from a Kubernetes pod label set by Argo Workflows; the \texttt{container="main"} filter excludes Argo's auxiliary \texttt{wait} and \texttt{init} containers whose logs are not relevant to the user.

The \texttt{env-vars/} module manages application secrets through a three-layer architecture. The NestJS service writes to and reads from HashiCorp Vault's KV v2 engine at the path \texttt{secret/data/projects/<tenantId>} using the Kubernetes Auth method, which authenticates the NestJS pod's ServiceAccount token against Vault without any static credential. Secret values are never stored in PostgreSQL — only metadata (key names, types, visibility, timestamps) is persisted relationally. After each successful Vault write, the service annotates the tenant's \texttt{ExternalSecret} resource with a timestamp trigger, causing the External Secrets Operator to immediately reconcile the Kubernetes Secret in the tenant namespace rather than waiting for the scheduled sync interval.

The \texttt{metrics/} module exposes two authenticated endpoints. \texttt{GET /metrics/dashboard} returns a per-user summary built from parallel Prometheus instant queries (request throughput, error rate, average response time, bandwidth, storage, CPU by project) combined with Prisma queries (active deployment count, failed builds, deployment history). \texttt{GET /metrics/projects/:id} returns five time-series for a given project (CPU, memory, response time, throughput, bandwidth) from Prometheus range queries, with resolution automatically adapted to the requested time window.

\subsection*{4.2 Go Build Worker}

The Go build worker is a single-purpose binary designed for execution as a Kubernetes Job within Argo Workflows. Its pipeline executes four stages in sequence within a 20-minute global timeout enforced via Go context. First, \texttt{git.go} performs a shallow clone (\texttt{--depth 1}) of the source repository, minimising data transfer for repositories with long histories. Second, \texttt{railpack.go} invokes \texttt{railpack prepare} against the cloned directory, producing \texttt{railpack-plan.json} and \texttt{railpack-info.json}. Railpack detects the application framework through static analysis of the repository file system — inspecting manifest files such as \texttt{package.json}, \texttt{go.mod}, \texttt{requirements.txt}, and \texttt{artisan} to identify the runtime, framework, and dependency manager. The preparation step is retried up to three times with clean state between attempts to absorb transient failures. Third, the generated JSON plan is patched by \texttt{patchRailpackPlanRunAsUser()}: the \texttt{deploy.user} field is set to \texttt{"1000"}, a \texttt{chown -R 1000:1000 /app} command is injected into the build steps, and \texttt{NODE\_OPTIONS=--max-old-space-size=<limit>} is appended to any Node.js start command to prevent V8 heap exhaustion under container memory limits. The memory limit value is read dynamically from \texttt{MEMORY\_LIMIT\_MB}, populated by Kubernetes resource field injection from the Job's resource limits. Fourth, \texttt{buildkit.go} connects to the Buildkitd daemon over gRPC, submits a \texttt{gateway.v0} frontend solve request using the Railpack frontend image as interpreter, streams build progress vertex-by-vertex to the Zap structured logger (collected by Promtail and indexed in Loki), and pushes the resulting OCI image to the internal registry.

\subsection*{4.3 Kubernetes Infrastructure}

The \texttt{tenant-chart} Helm chart provisions all Kubernetes resources required to run and expose a tenant application in a single atomic operation. For each deployment, it creates: a dedicated \texttt{Namespace} labelled with tenant identity; a \texttt{ServiceAccount} with \texttt{automountServiceAccountToken: false}; a \texttt{Deployment} with \texttt{runAsNonRoot: true}, \texttt{runAsUser: 1000}, \texttt{topologySpreadConstraints} for cross-node distribution, and HTTP readiness/liveness probes; a \texttt{ClusterIP Service}; a Traefik \texttt{IngressRoute} matching \texttt{Host(<subdomain>.hostifer.me)}; a \texttt{LimitRange} enforcing per-container resource envelopes; a \texttt{ResourceQuota} capping namespace aggregate consumption and explicitly setting \texttt{persistentvolumeclaims: 0} for stateless applications; and four \texttt{NetworkPolicies} implementing a strict allowlist model: default deny all, allow ingress from Traefik namespace only, allow DNS egress on UDP/TCP 53, and allow egress to public internet excluding RFC1918 private ranges. Resource allocation is plan-dependent, ranging from 100m CPU / 128Mi RAM (Free) to 1 CPU / 1Gi RAM (Enterprise), with HPA enabled from Standard plan upward and VPA available in \texttt{recommendation} mode (Standard/Pro) or \texttt{auto} mode (Enterprise).

Two Argo Workflows WorkflowTemplates orchestrate the deployment and deletion lifecycles. The \texttt{hostifer-deploy} template executes three sequential steps: \textit{Build} (running the Go builder container), \textit{Provision} (running \texttt{helm install} with the tenant chart from GHCR OCI registry with \texttt{--wait} and \texttt{--create-namespace}), and \textit{DNS} (running the \texttt{create-dns-job} Go binary to create a Cloudflare DNS record for the subdomain). The \texttt{hostifer-delete} template executes three parallel steps: Helm uninstall, DNS record deletion, and OCI registry image removal via the Docker Distribution HTTP API. An \texttt{onExit} handler triggers a Next.js reconciliation job to update the frontend state after deletion. RBAC is managed through two ClusterRoles: \texttt{argo-provisioner} (granting cross-namespace resource creation for Helm operations) and \texttt{argo-workflow-controller} (granting Argo CRD management), both bound to the \texttt{argo-workflow-sa} ServiceAccount.

TLS certificate lifecycle is automated by cert-manager configured with a \texttt{ClusterIssuer} using the Let's Encrypt ACME protocol and Cloudflare DNS-01 challenge. A wildcard certificate \texttt{*.hostifer.me} covers all tenant subdomains without requiring a per-application certificate issuance, with automatic renewal triggered 30 days before expiry. Observability is provided by Promtail (DaemonSet) collecting pod logs and attaching \texttt{build\_id}, \texttt{tenant\_id}, \texttt{container}, and \texttt{log\_type} labels via Kubernetes service discovery relabeling, forwarding to Loki for indexing and retention; and by Prometheus scraping Traefik, cAdvisor, kubelet, and NestJS \texttt{/metrics} endpoints via ServiceMonitors, with metrics exposed to the HPA via the Prometheus Adapter for custom metric autoscaling.

% ── 5. Testing ────────────────────────────────────────────────
\section*{5. Testing and Validation}

\subsection*{5.1 Functional Testing}

Functional validation was conducted by deploying reference applications for eight frameworks supported by Railpack: Go, Express.js, Flask, NestJS, Django, Laravel, Next.js, and Spring Boot. Each framework was tested with five consecutive deployment runs on identical source code. Mean deployment times ranged from 1 minute 50 seconds (Go binary compilation) to 4 minutes 50 seconds (Spring Boot with Maven dependency resolution), all within the 5-minute BNF-01 target. All 40 deployment runs completed successfully, with two transient Spring Boot timeouts automatically resolved by the \texttt{GenerateBuildWithRetry()} mechanism. Secret injection from Vault via ESO was validated for all frameworks, with secrets available in tenant pods within 30 seconds of the Provision step completion and no sensitive values appearing in Loki-indexed build logs.

\subsection*{5.2 Network Isolation Testing}

Inter-tenant network isolation was validated by attempting direct TCP connections between pods in different tenant namespaces using \texttt{netcat} and \texttt{curl} from within running containers. All 12 connection attempts were rejected by the Kubernetes network plugin enforcing the \texttt{default-deny-all} NetworkPolicy, producing immediate connection refused responses. Ingress from the Traefik namespace was confirmed functional across all tested tenant namespaces. DNS resolution (UDP 53) was confirmed available from all tenant pods. Egress to RFC1918 addresses was confirmed blocked, while egress to public internet endpoints was confirmed permitted, validating the four-policy model implemented in the tenant chart.

\subsection*{5.3 Performance and Autoscaling}

Performance tests measured three concurrent build scenarios (1, 3, and 5 simultaneous builds) against the shared Buildkitd daemon. Buildkitd memory consumption scaled from approximately 450 MiB under a single build to approximately 1.8 GiB under five concurrent builds, remaining within the 2 GiB limit configured for the Buildkitd pod. SSE log latency remained below 2.5 seconds for all tested scenarios. HPA validation under simulated CPU load showed replica convergence within 25 to 45 seconds of threshold crossing, approaching the BNF-09 target of 30 seconds. Cluster Autoscaler node provisioning time averaged between 2 minutes 30 seconds and 4 minutes, slightly exceeding the BNF-11 target of 3 minutes under high initial cluster load.

\subsection*{5.4 Regulatory Compliance}

The compliance checklist verified eight categories of regulatory and technical requirements. Data residency was confirmed: all platform components — PostgreSQL, Vault, Loki, the OCI registry, and tenant application pods — operate exclusively within the Alibaba Cloud Algeria infrastructure. TLS encryption was confirmed active on all user-facing endpoints. Secret non-exposure was verified by confirming that no \texttt{env-vars} secret values appear in Loki logs, PostgreSQL records, or Kubernetes resource manifests. Audit log creation was verified for all sensitive operations (secret creation, modification, deletion, and reveal). HMAC webhook signature validation was confirmed rejecting requests with invalid signatures. Non-root pod execution was confirmed for all tenant workloads. Token automount disabling was confirmed on all tenant ServiceAccounts.

% ── 6. Results and Perspectives ──────────────────────────────
\section*{6. Results, Challenges, and Perspectives}

\subsection*{6.1 Achieved Results}

Of the 20 functional requirements and 14 non-functional requirements identified in Chapter 2, 17 functional requirements were fully implemented (85\%), 2 were partially implemented (rollback UI and administration interface), and 1 was deferred (team member management). Of the non-functional requirements, 10 were fully satisfied, 4 were partially satisfied (SLA measurement, HPA convergence time, Cluster Autoscaler latency, and VPA auto mode in production), and none were deferred. The platform successfully demonstrates that a native Kubernetes PaaS with developer experience comparable to international standards can be built, hosted, and operated within Algeria, in compliance with the applicable legal framework, and priced accessibly in DZD.

\subsection*{6.2 Technical Challenges}

Four significant technical challenges were encountered and resolved during implementation. The Railpack binary on Alpine Linux images presents a musl/glibc compatibility issue requiring a custom base image build or the use of a Debian-based runtime, resolved by adopting the Railpack-provided runtime base image. Loki \texttt{query\_range} log parsing required careful handling of nanosecond timestamps and duplicate line deduplication across polling intervals, resolved by tracking delivered line signatures in the async generator state. NextAuth.js refresh token handling in the JWT callback required precise error state propagation to prevent silent session invalidation, resolved by a dedicated \texttt{error} field in the token object. Argo Workflows RBAC for cross-namespace Helm provisioning required the construction of a custom \texttt{ClusterRole} combining namespace creation, Helm secret management, and Traefik CRD permissions, as no built-in role provides this combination.

\subsection*{6.3 Future Perspectives}

Four priority directions have been identified for future development. Cold-start elimination using KEDA (Kubernetes Event-Driven Autoscaler) would allow tenant applications on the Free plan to scale to zero during inactivity and resume on demand, significantly reducing infrastructure cost for low-traffic applications. Managed database provisioning would allow developers to attach PostgreSQL, MySQL, or Redis instances to their applications directly from the platform, eliminating the need to connect to external database services. S3-compatible object storage, exposed via MinIO deployed on the cluster, would provide stateful application support within the existing infrastructure. GitHub organisation support would allow teams to connect repositories from GitHub organisations rather than individual accounts, enabling collaborative project management within the current access model.

% ── 7. Conclusion ─────────────────────────────────────────────
\section*{7. Conclusion}

Hostifer demonstrates that the gap between the developer experience offered by leading international PaaS platforms and the practical accessibility, legal compliance, and affordability required by the Algerian digital ecosystem can be bridged through a carefully designed cloud-native architecture. By combining Kubernetes-native multi-tenant isolation, an automated build pipeline based on Railpack and Buildkit, Vault-backed secret management, automated TLS via cert-manager and Let's Encrypt, and real-time observability through the Promtail-Loki-Prometheus stack — all hosted on Algerian infrastructure — the platform addresses simultaneously the six structural pain points identified in the market analysis: USD-denominated pricing, regulatory non-compliance, absence of automated deployment, lack of robust multi-tenant isolation, geographic latency, and missing integrated observability.

The platform's contribution operates on three axes. Technically, it demonstrates the feasibility of a production-grade Kubernetes multi-tenant PaaS built by a small team using exclusively open-source components. Economically, it provides a DZD-native, free-tier-inclusive pricing model that removes the financial barrier to professional cloud deployment for Algerian developers. Societally, it contributes to digital sovereignty by providing an infrastructure layer that ensures personal data of Algerian users remains on Algerian territory, in compliance with the legal framework, without sacrificing developer experience. As the Algerian regulatory environment continues to evolve toward stronger data residency enforcement, platforms such as Hostifer represent not merely a convenience but an infrastructural necessity for the growth of a compliant and sovereign Algerian digital economy.
